\documentclass[12pt,a4paper]{article}

% -----------------------------------------------------------------------
% Packages
% -----------------------------------------------------------------------
\usepackage[utf8]{inputenc}
\usepackage[T1]{fontenc}
\usepackage{geometry}
\usepackage{graphicx}
\usepackage{hyperref}
\usepackage{booktabs}
\usepackage{tabularx}
\usepackage{longtable}
\usepackage{array}
\usepackage{enumitem}
\usepackage{fancyhdr}
\usepackage[table]{xcolor}
\usepackage{titlesec}
\usepackage{datetime2}
\usepackage{amsmath}

\geometry{margin=2.5cm}

% -----------------------------------------------------------------------
% Colours & formatting
% -----------------------------------------------------------------------
\definecolor{sscblue}{RGB}{0,70,127}
\definecolor{lightgray}{RGB}{230,230,230}

\titleformat{\section}{\large\bfseries\color{sscblue}}{{\thesection}}{1em}{}[\titlerule]
\titleformat{\subsection}{\normalsize\bfseries\color{sscblue}}{{\thesubsection}}{1em}{}

% -----------------------------------------------------------------------
% Header / footer
% -----------------------------------------------------------------------
\pagestyle{fancy}
\fancyhf{}
\rhead{\small SSC\_D-M001 \textbar{} Rev~1.0}
\lhead{\small Project Management Plan}
\rfoot{\small Page \thepage}
\lfoot{\small \textcopyright{} SSC -- CONFIDENTIAL}

% -----------------------------------------------------------------------
% Hyperref setup
% -----------------------------------------------------------------------
\hypersetup{
  colorlinks = true,
  linkcolor  = sscblue,
  urlcolor   = sscblue,
  pdftitle   = {SSC\_D-M001 -- Project Management Plan},
  pdfauthor  = {SSC Team},
}

% -----------------------------------------------------------------------
% Document metadata (edit here for each release)
% -----------------------------------------------------------------------
\newcommand{\docid}{SSC\_D-M001}
\newcommand{\doctitle}{Project Management Plan}
\newcommand{\docrev}{1.0}
\newcommand{\docdate}{\DTMdate{2026-04-03}}
\newcommand{\docstatus}{Draft}
\newcommand{\docauthor}{SSC Team}

% =======================================================================
\begin{document}
% =======================================================================

% -----------------------------------------------------------------------
% Title page
% -----------------------------------------------------------------------
\begin{titlepage}
  \centering
  \vspace*{2cm}

  {\Huge\bfseries\color{sscblue} \doctitle \par}
  \vspace{0.5cm}
  {\large Document ID: \texttt{\docid} \par}
  \vspace{2cm}

  \begin{tabular}{ll}
    \textbf{Revision:} & \docrev \\[4pt]
    \textbf{Date:}     & \docdate \\[4pt]
    \textbf{Status:}   & \docstatus \\[4pt]
    \textbf{Author:}   & \docauthor \\
  \end{tabular}

  \vfill
  \rule{\linewidth}{0.4pt}\\[4pt]
  {\small This document is subject to change without notice.\\
  All rights reserved -- SSC Project.}
\end{titlepage}

% -----------------------------------------------------------------------
% Revision history
% -----------------------------------------------------------------------
\newpage
\section*{Document Revision History}
\begin{tabularx}{\textwidth}{|c|c|X|l|}
  \hline
  \rowcolor{lightgray}
  \textbf{Rev} & \textbf{Date} & \textbf{Description} & \textbf{Author} \\
  \hline
  1.0 & 2026-04-03 & Initial release & SSC Team \\
  \hline
\end{tabularx}

\newpage
\tableofcontents
\newpage

% -----------------------------------------------------------------------
% 1. Introduction
% -----------------------------------------------------------------------
\section{Introduction}

\subsection{Purpose}
This document defines the Project Management Plan (PMP) for the SSC project.
It describes the governance structure, planning methodology, risk management
approach, and quality assurance processes adopted throughout the project
lifecycle.

\subsection{Scope}
This plan applies to all activities within the SSC project from initiation
through to closure. It is applicable to all project team members, stakeholders,
and subcontractors operating under the SSC programme.

\subsection{Document Naming Convention}
Documents produced under this project follow the naming convention:

\begin{center}
  \texttt{SSC\_D-<Category><NNN>}
\end{center}

\noindent where \texttt{<Category>} is one of:

\begin{description}[labelindent=1cm]
  \item[\texttt{C}] Configuration \& Interface Control
  \item[\texttt{T}] Technical Specifications \& Design
  \item[\texttt{M}] Management \& Planning
  \item[\texttt{D}] Data, Analysis \& Reports
\end{description}

\noindent and \texttt{<NNN>} is a zero-padded three-digit sequence number.

\subsection{Applicable Documents}
\begin{enumerate}
  \item SSC\_D-T001 -- System Requirements Specification
  \item SSC\_D-C001 -- Interface Control Document
  \item SSC\_D-D001 -- Verification \& Validation Report
\end{enumerate}

% -----------------------------------------------------------------------
% 2. Project Overview
% -----------------------------------------------------------------------
\section{Project Overview}

\subsection{Objectives}
The primary objectives of the SSC project are:
\begin{itemize}
  \item Deliver a fully verified system in accordance with the approved
        requirements baseline.
  \item Maintain schedule, cost, and quality within agreed tolerances.
  \item Ensure traceability from requirements through design to verification.
\end{itemize}

\subsection{Deliverables}
Table~\ref{tab:deliverables} lists the key deliverables.

\begin{longtable}{|c|l|l|c|}
  \hline
  \rowcolor{lightgray}
  \textbf{ID} & \textbf{Deliverable} & \textbf{Document} & \textbf{Due} \\
  \hline
  \endfirsthead
  \hline
  \rowcolor{lightgray}
  \textbf{ID} & \textbf{Deliverable} & \textbf{Document} & \textbf{Due} \\
  \hline
  \endhead
  D-01 & Project Management Plan   & SSC\_D-M001 & Q1 2026 \\
  D-02 & System Requirements Spec  & SSC\_D-T001 & Q1 2026 \\
  D-03 & Preliminary Design Review & SSC\_D-T002 & Q2 2026 \\
  D-04 & Critical Design Review    & SSC\_D-T003 & Q3 2026 \\
  D-05 & V\&V Report               & SSC\_D-D001 & Q4 2026 \\
  \hline
  \caption{Project Deliverables}
  \label{tab:deliverables}
\end{longtable}

% -----------------------------------------------------------------------
% 3. Organisation & Roles
% -----------------------------------------------------------------------
\section{Organisation and Roles}

\subsection{Project Governance}
The project operates under a standard governance framework consisting of:

\begin{itemize}
  \item \textbf{Project Board} -- strategic oversight and go/no-go decisions.
  \item \textbf{Project Manager} -- day-to-day management and reporting.
  \item \textbf{Technical Lead} -- technical authority and design decisions.
  \item \textbf{Quality Manager} -- process compliance and audit.
\end{itemize}

\subsection{Responsibilities Matrix (RACI)}
\begin{tabularx}{\textwidth}{|X|c|c|c|c|}
  \hline
  \rowcolor{lightgray}
  \textbf{Activity} & \textbf{PM} & \textbf{TL} & \textbf{QM} & \textbf{Board} \\
  \hline
  Schedule management   & R & C & I & A \\
  Technical decisions   & I & R & I & A \\
  Risk management       & R & C & C & A \\
  Quality audits        & I & I & R & A \\
  Budget control        & R & I & I & A \\
  \hline
\end{tabularx}

\vspace{6pt}
\noindent\textit{R = Responsible, A = Accountable, C = Consulted, I = Informed}

% -----------------------------------------------------------------------
% 4. Risk Management
% -----------------------------------------------------------------------
\section{Risk Management}

\subsection{Risk Assessment Process}
Risks are assessed using a $5 \times 5$ likelihood--impact matrix.
Risk exposure is calculated as:

\begin{equation}
  E = L \times I
\end{equation}

\noindent where $L \in [1,5]$ is the likelihood score and $I \in [1,5]$ is
the impact score.  Risks with $E \geq 12$ are escalated to the Project Board.

\subsection{Risk Register (excerpt)}
\begin{tabularx}{\textwidth}{|c|X|c|c|c|X|}
  \hline
  \rowcolor{lightgray}
  \textbf{ID} & \textbf{Risk} & \textbf{L} & \textbf{I} & \textbf{E} & \textbf{Mitigation} \\
  \hline
  R-01 & Key personnel unavailability & 3 & 4 & 12 & Cross-train team members \\
  R-02 & Requirements change late     & 2 & 5 & 10 & Freeze baseline at PDR \\
  R-03 & CI pipeline failure          & 2 & 3 &  6 & Maintain local build fallback \\
  \hline
\end{tabularx}

% -----------------------------------------------------------------------
% 5. Quality Assurance
% -----------------------------------------------------------------------
\section{Quality Assurance}

\subsection{Document Control}
All project documents are version-controlled in the project repository.
PDF artefacts are generated automatically via the Continuous Integration
pipeline described in Section~\ref{sec:ci}.

\subsection{Continuous Integration}\label{sec:ci}
The project uses GitHub Actions to automate the following tasks:

\begin{itemize}
  \item Compile all \LaTeX{} sources to PDF on every pull request.
  \item Publish the compiled PDF as a build artefact for review.
  \item Attach the signed PDF to GitHub Releases at release time.
\end{itemize}

% -----------------------------------------------------------------------
% 6. Glossary
% -----------------------------------------------------------------------
\section{Glossary}
\begin{tabularx}{\textwidth}{|l|X|}
  \hline
  \rowcolor{lightgray}
  \textbf{Term} & \textbf{Definition} \\
  \hline
  CI   & Continuous Integration \\
  PMP  & Project Management Plan \\
  PDR  & Preliminary Design Review \\
  CDR  & Critical Design Review \\
  RACI & Responsible, Accountable, Consulted, Informed \\
  SSC  & System / Sub-System Controller (project code) \\
  V\&V & Verification and Validation \\
  \hline
\end{tabularx}

\end{document}
